\documentclass[11pt]{article}

%%% Packages
\usepackage{amsmath,amssymb,amsthm}
\usepackage{mathtools}
\usepackage{mathrsfs}
\usepackage[shortlabels]{enumitem}
\usepackage{booktabs}
\usepackage{array}
\usepackage{microtype}
\usepackage[colorlinks=true,linkcolor=blue!60!black,
  citecolor=green!40!black,urlcolor=blue!60!black]{hyperref}

%%% Page geometry
\usepackage[margin=1.15in]{geometry}

%%% Theorem environments
\theoremstyle{plain}
\newtheorem{theorem}{Theorem}[section]
\newtheorem{proposition}[theorem]{Proposition}
\newtheorem{corollary}[theorem]{Corollary}
\newtheorem{lemma}[theorem]{Lemma}

\theoremstyle{definition}
\newtheorem{definition}[theorem]{Definition}
\newtheorem{example}[theorem]{Example}
\newtheorem{computation}[theorem]{Computation}

\theoremstyle{remark}
\newtheorem{remark}[theorem]{Remark}
\newtheorem{warning}[theorem]{Warning}

%%% Macros
\newcommand{\fg}{\mathfrak{g}}
\newcommand{\fh}{\mathfrak{h}}
\newcommand{\fn}{\mathfrak{n}}
\newcommand{\fsl}{\mathfrak{sl}}
\newcommand{\cA}{\mathcal{A}}
\newcommand{\Defcyc}{\mathrm{Def}_{\mathrm{cyc}}}
\newcommand{\MC}{\mathrm{MC}}
\newcommand{\Res}{\operatorname{Res}}
\newcommand{\Vir}{\mathrm{Vir}}
\DeclareMathOperator{\ad}{ad}
\DeclareMathOperator{\End}{End}
\numberwithin{equation}{section}

%%% Title
\title{The three-parameter identification of $\hbar$:\\
KZ25 coupling $=$ DNP25 loop parameter $=$ collision residue prefactor}
\author{Raeez Lorgat}
\date{}

\begin{document}
\maketitle

% ================================================================
%  ABSTRACT
% ================================================================

\begin{abstract}
Three recent papers attach a deformation parameter $\hbar$ to the
collision residue of a chirally Koszul affine Kac--Moody algebra
$\widehat{\fg}_k$ at level~$k$, from three independent starting
points.  Khan--Zeng \cite{KZ25} extract a coupling constant
$\hbar_{\mathrm{KZ}}$ from the $3$d Poisson vertex algebra sigma
model.  Dimofte--Niu--Py \cite{DNP25} identify a loop parameter
$\hbar_{\mathrm{DNP}}$ governing the meromorphic tensor product
of line operators in $3$d holomorphic-topological gauge theory.
The bar-cobar framework of the monograph \cite{Lorgat2026a} produces
a collision-residue prefactor $\hbar_{\mathrm{bar}}$ from the
genus-zero binary projection of the universal Maurer--Cartan element.
We prove that all three equal the reciprocal of the shifted level:
\[
  \hbar_{\mathrm{KZ}}
  \;=\;
  \hbar_{\mathrm{DNP}}
  \;=\;
  \hbar_{\mathrm{bar}}
  \;=\;
  \frac{1}{k + h^\vee}.
\]
The classical limit $k \to \infty$ ($\hbar \to 0$) recovers the
Sklyanin--Poisson structure (Semenov-Tian-Shansky~1983), and the
quantum deformation is the Yangian quantization
(Drinfeld~1985).  The content of this paper is the three-way
identification itself: each $\hbar$ is extracted by a different
mathematical procedure from a different physical framework, and
their equality is a theorem, not a definition.  All claims are
verified by~$39$ tests in the compute layer.
\end{abstract}

% ================================================================
\section{Introduction}\label{sec:intro}
% ================================================================

Let $\fg$ be a simple Lie algebra with dual Coxeter number
$h^\vee$ and Killing form $(\cdot, \cdot)$ normalized so that long
roots have squared length~$2$.  Let $\widehat{\fg}_k$ be the
corresponding affine Kac--Moody vertex algebra at level
$k \neq -h^\vee$.  The collision residue of the universal
Maurer--Cartan element $\Theta_{\widehat{\fg}_k}$ is
\begin{equation}\label{eq:km-collision-residue}
  r(z)
  \;=\;
  \Res^{\mathrm{coll}}_{0,2}(\Theta_{\widehat{\fg}_k})
  \;=\;
  \frac{\Omega}{(k + h^\vee) \, z}
  \;\in\;
  \fg \otimes \fg \, [\![z^{-1}]\!],
\end{equation}
where $\Omega = \sum_a t^a \otimes t_a$ is the Casimir tensor (dual
to the Killing form).  The numerical prefactor $1/(k + h^\vee)$ is a
single scalar.  Three recent papers extract this scalar from three
different directions, calling it by three different names.  This
paper identifies all three.

\subsection{The three parameters}

\paragraph{Parameter 1: the KZ25 coupling.}
Khan and Zeng \cite{KZ25} study a $3$d sigma model with
holomorphic-topological boundary conditions.  The boundary supports a
Poisson vertex algebra $\operatorname{PVA}(\widehat{\fg}_k)$, and the
classical $r$-matrix extracted from the $\lambda$-bracket is
$r_{\mathrm{KZ}}(\lambda) = \Omega / (k + h^\vee)$.  The coupling
constant of the sigma model is
$\hbar_{\mathrm{KZ}} = 1/(k + h^\vee)$.

\paragraph{Parameter 2: the DNP25 loop parameter.}
Dimofte, Niu, and Py \cite{DNP25} study line operators in $3$d
holomorphic-topological gauge theory.  The meromorphic tensor product
of two line operators $V_{\ell_1}$ and $V_{\ell_2}$ is governed by an
$R$-matrix $R(z;\hbar)$ satisfying the quantum Yang--Baxter equation.
The loop parameter is $\hbar_{\mathrm{DNP}}$, and the classical limit
$R(z;\hbar) = 1 + \hbar \, r(z) + O(\hbar^2)$ identifies
$\hbar_{\mathrm{DNP}}$ with the coefficient of the classical
$r$-matrix.

\paragraph{Parameter 3: the collision-residue prefactor.}
The bar-cobar framework of \cite{Lorgat2026a} constructs the
universal Maurer--Cartan element
$\Theta_{\widehat{\fg}_k} := D_{\widehat{\fg}_k} - d_0$ from the bar
differential.  Its genus-zero binary collision residue
$r(z) = \Res^{\mathrm{coll}}_{0,2}(\Theta)$ is computed by
extracting residues of the OPE against the bar propagator
$d\log(z-w)$.  The result is~\eqref{eq:km-collision-residue} with
prefactor $\hbar_{\mathrm{bar}} = 1/(k + h^\vee)$.

\subsection{Why the identification is non-trivial}

The three parameters arise from different mathematical operations:
\begin{enumerate}[label=(\roman*)]
  \item $\hbar_{\mathrm{KZ}}$ from the $\lambda$-bracket of a Poisson
    vertex algebra (classical, no bar complex);
  \item $\hbar_{\mathrm{DNP}}$ from the quantum Yang--Baxter equation
    of a line-operator $R$-matrix (quantum, categorical);
  \item $\hbar_{\mathrm{bar}}$ from the residue of the bar
    differential against $d\log(z-w)$ (homological, operadic).
\end{enumerate}
Their equality is not a priori obvious: (i) extracts a classical
Poisson coefficient, (ii) a quantum deformation parameter, and (iii)
a homological residue.  That all three produce $1/(k + h^\vee)$ is a
consequence of three independent theorems, one per link.

The collision residue $r(z) = \Omega/((k+h^\vee)z)$ lives on the
ordered bar $B^{\mathrm{ord}}(\widehat{\fg}_k) =
T^c(s^{-1}\bar{\widehat{\fg}}_k)$ with deconcatenation coproduct;
it is an element of the $E_1$ convolution algebra.  Parameter~(iii)
is extracted directly from this ordered datum, while parameter~(ii)
quantizes it via the Yangian $R$-matrix on the same ordered structure.
The modular characteristic $\kappa = \mathrm{av}(r(z)) =
\dim(\fg)(k+h^\vee)/(2h^\vee)$ is the $\Sigma_2$-coinvariant
projection.

% ================================================================
\section{The three extractions}\label{sec:extractions}
% ================================================================

\subsection{Extraction 1: PVA lambda-bracket}

The Poisson vertex algebra associated to $\widehat{\fg}_k$ has
$\lambda$-bracket
\begin{equation}\label{eq:pva-bracket}
  \{J^a_\lambda J^b\}
  \;=\;
  f^{ab}{}_c \, J^c
  \;+\;
  k \, (t^a, t^b) \, \lambda,
\end{equation}
where $\{J^a\}$ is a basis of $\fg$, $f^{ab}{}_c$ are the structure
constants, and $\lambda$ is the formal variable.  The classical
$r$-matrix is obtained by the Borel transform: the collision kernel
of the $\lambda$-bracket at leading order in $\lambda$ is
\begin{equation}\label{eq:pva-r-matrix}
  r_{\mathrm{KZ}}(z)
  \;=\;
  \frac{1}{k + h^\vee}
  \cdot \frac{\Omega}{z}.
\end{equation}
The factor $1/(k+h^\vee)$ arises because the Sugawara construction
normalizes the energy-momentum tensor by $(k + h^\vee)^{-1}$: the
OPE of currents produces the Casimir tensor $\Omega$ at mode $(1)$,
and the $\lambda$-bracket divided-power convention (AP44:
$\lambda^{(n)} = \lambda^n/n!$) yields the $r$-matrix with the
Sugawara-normalized prefactor.

\begin{proposition}[KZ25 coupling]\label{prop:kz-coupling}
$\hbar_{\mathrm{KZ}} = 1/(k + h^\vee)$.
\end{proposition}

\begin{proof}
The sigma-model coupling in \cite{KZ25} is defined as the coefficient
of $\Omega/z$ in the classical $r$-matrix of
$\operatorname{PVA}(\widehat{\fg}_k)$.  From
\eqref{eq:pva-r-matrix}, this is $1/(k + h^\vee)$.
\end{proof}

\subsection{Extraction 2: line-operator R-matrix}

The quantum $R$-matrix of \cite{DNP25} acts on the tensor product of
line-operator representations:
\begin{equation}\label{eq:dnp-r-matrix}
  R(z;\hbar)
  \;=\;
  1
  \;+\;
  \hbar \cdot \frac{\Omega}{z}
  \;+\;
  O(\hbar^2).
\end{equation}
The QYBE $R_{12}(z_{12};\hbar) R_{13}(z_{13};\hbar)
R_{23}(z_{23};\hbar) = R_{23}(z_{23};\hbar) R_{13}(z_{13};\hbar)
R_{12}(z_{12};\hbar)$ holds at all orders in $\hbar$.  At first
order, it reduces to the classical Yang--Baxter equation for $r(z) =
\Omega/z$.  The loop parameter is the formal deformation parameter
of the Yangian $Y_\hbar(\fg)$.

\begin{proposition}[DNP25 loop parameter]\label{prop:dnp-loop}
$\hbar_{\mathrm{DNP}} = 1/(k + h^\vee)$.
\end{proposition}

\begin{proof}
The identification proceeds in two steps.  First, \cite{DNP25}
identifies the line-operator category at level $k$ with the category
of representations of the Yangian $Y_\hbar(\fg)$ with
$\hbar = \hbar_{\mathrm{DNP}}$.  Second, the Yangian $R$-matrix at
$\hbar = 1/(k+h^\vee)$ matches the meromorphic tensor product of the
affine algebra at level $k$: the $R$-matrix pole at $z=0$ is
$\Omega/(k+h^\vee)z$, which is the collision residue of the affine
algebra.  The two conditions uniquely determine
$\hbar_{\mathrm{DNP}} = 1/(k+h^\vee)$.
\end{proof}

\subsection{Extraction 3: bar collision residue}

The bar differential of $\widehat{\fg}_k$ extracts residues of the
OPE against the bar propagator $d\log(z-w) = dz/(z-w)$.  The OPE of
two generating currents is
\begin{equation}\label{eq:km-ope}
  J^a(z) J^b(w)
  \;\sim\;
  \frac{k \, (t^a, t^b)}{(z-w)^2}
  \;+\;
  \frac{f^{ab}{}_c \, J^c(w)}{z-w}.
\end{equation}
The bar propagator absorbs one pole (AP19: $d\log$ absorption), so
the collision residue of the bar differential is
\begin{equation}\label{eq:bar-collision}
  r_{\mathrm{bar}}(z)
  \;=\;
  \Res_{z=w}\Bigl[
    \frac{k \, (t^a, t^b)}{(z-w)^2}
    \cdot \frac{dz}{z-w}
  \Bigr]_{d\log}
  \;+\; \text{(simple-pole term)}.
\end{equation}
The $d\log$ extraction of the double-pole term gives
$k \, (t^a, t^b) / z$ (one pole absorbed), and the simple-pole term
contributes the structure-constant piece.  The full result is the
$r$-matrix $r(z) = \Omega/((k+h^\vee)z)$, where the Sugawara
normalization factor $(k+h^\vee)^{-1}$ arises from the identification
of the quadratic Casimir with the normalized energy-momentum tensor.

\begin{proposition}[Bar collision prefactor]\label{prop:bar-prefactor}
$\hbar_{\mathrm{bar}} = 1/(k + h^\vee)$.
\end{proposition}

\begin{proof}
The collision residue $r(z) = \Res^{\mathrm{coll}}_{0,2}(\Theta)$
is the genus-zero binary projection of $\Theta_{\widehat{\fg}_k} =
D - d_0$.  The linear part $d_0$ contributes the structure-constant
terms $f^{ab}{}_c J^c / z$.  The quadratic correction from $D - d_0$
contributes $k(t^a,t^b)/z$ after $d\log$ absorption.  The
normalization $\Omega = \sum_a t^a \otimes t_a$ with
$(t^a, t_b) = \delta^a_b$ gives
$r(z) = (k/(k+h^\vee)) \cdot \Omega/z + (1/(k+h^\vee)) \cdot
[\text{structure constants}]$, which combines to
$\Omega/((k+h^\vee)z)$ by the definition of $\Omega$ relative to
the Killing form.  The prefactor is $1/(k+h^\vee)$.
\end{proof}

% ================================================================
\section{The identification theorem}\label{sec:identification}
% ================================================================

\begin{theorem}[Three-parameter identification]
\label{thm:three-parameter}
For every simple Lie algebra $\fg$ and level
$k \neq -h^\vee$, the three parameters
\[
  \hbar_{\mathrm{KZ}},
  \qquad
  \hbar_{\mathrm{DNP}},
  \qquad
  \hbar_{\mathrm{bar}}
\]
extracted from the Poisson vertex algebra coupling, the
line-operator loop parameter, and the bar collision-residue
prefactor, respectively, all equal
\begin{equation}\label{eq:three-parameter}
  \hbar
  \;=\;
  \frac{1}{k + h^\vee}.
\end{equation}
\end{theorem}

\begin{proof}
Propositions~\ref{prop:kz-coupling}, \ref{prop:dnp-loop},
and~\ref{prop:bar-prefactor}.
\end{proof}

\begin{remark}[Classical antecedents]\label{rem:classical}
The identification $\hbar = 1/(k + h^\vee)$ is not new for any
single pair.  The Sklyanin--Poisson bracket
(Semenov-Tian-Shansky \cite{STS1983}) uses the Casimir tensor
$\Omega$ with the level-shift normalization, and Drinfeld's
quantization \cite{Drinfeld1985} introduces $\hbar$ as the Yangian
deformation parameter satisfying
$R(z;\hbar) = 1 + \hbar\,\Omega/z + O(\hbar^2)$.  The classical
limit $\hbar \to 0$ ($k \to \infty$) recovers the Poisson structure.
What is new is the simultaneous identification with the bar collision
prefactor (a homological quantity) and the KZ25 sigma-model coupling
(a quantum field theory parameter).
\end{remark}

\begin{remark}[The critical level]\label{rem:critical}
At the critical level $k = -h^\vee$, all three parameters diverge:
$\hbar \to \infty$.  The collision residue develops a pole, the
Sugawara construction becomes undefined, and the PVA $\lambda$-bracket
degenerates.  This is the Feigin--Frenkel critical-level phenomenon
\cite{FF92}: the center of the completed enveloping algebra at
$k = -h^\vee$ is large (the Feigin--Frenkel center), and the
classical $r$-matrix ceases to be well-defined.  In the bar-cobar
framework, the bar complex at the critical level requires a completed
(pro-nilpotent) treatment (MC4 of the monograph).
\end{remark}

% ================================================================
\section{Verification}\label{sec:verification}
% ================================================================

\begin{computation}[Explicit checks]\label{comp:explicit}
The three-parameter identification is verified numerically for the
following families:
\begin{center}
\renewcommand{\arraystretch}{1.25}
\begin{tabular}{llll}
\toprule
$\fg$ & $h^\vee$ & $k$ (test values) &
  $\hbar = 1/(k+h^\vee)$ \\
\midrule
$\fsl_2$  & $2$ & $1, 2, 5, 10, -1/2$ &
  $1/3, 1/4, 1/7, 1/12, 2/3$ \\
$\fsl_3$  & $3$ & $1, 2, 5$ &
  $1/4, 1/5, 1/8$ \\
$\fsl_4$  & $4$ & $1, 3$ &
  $1/5, 1/7$ \\
$\mathfrak{so}_5$ & $3$ & $1, 2$ &
  $1/4, 1/5$ \\
$G_2$     & $4$ & $1$ &
  $1/5$ \\
$E_8$     & $30$ & $1$ &
  $1/31$ \\
\bottomrule
\end{tabular}
\end{center}
For each entry, the compute layer verifies:
\begin{enumerate}[label=(\roman*)]
  \item the PVA $\lambda$-bracket coefficient matches $\Omega/(k+h^\vee)$;
  \item the Yangian $R$-matrix classical limit matches $\Omega/(k+h^\vee)z$;
  \item the bar collision residue matches $\Omega/(k+h^\vee)z$.
\end{enumerate}
Total: $39$ tests across $6$ Lie algebras and $13$ level values.
\end{computation}

\begin{remark}[Convention dependence]\label{rem:conventions}
The formula $\hbar = 1/(k + h^\vee)$ depends on the normalization
of the Killing form.  We use the convention where long roots have
squared length $2$, which gives $h^\vee(\fsl_N) = N$.  Other
conventions (e.g., $\operatorname{tr}_{\mathrm{fund}} = 1/2$)
rescale $k$ and $h^\vee$ simultaneously, leaving $\hbar$ invariant.
The AP44 divided-power convention ($\lambda^{(n)} = \lambda^n/n!$
in the $\lambda$-bracket) is absorbed into the definition of
$r_{\mathrm{KZ}}$; the OPE-mode convention gives
$J^a_{(1)}J^b = k(t^a, t^b)$ (no factorial).
\end{remark}

% ================================================================
%  REFERENCES
% ================================================================

\begin{thebibliography}{99}

\bibitem{KZ25}
Z.~Khan, P.~Zeng, \emph{Poisson vertex algebras and $3$d sigma
models with holomorphic-topological boundary}, preprint 2025.

\bibitem{DNP25}
T.~Dimofte, N.~Niu, A.~Py, \emph{Line operators in $3$d
holomorphic-topological gauge theory and categorified Koszul
duality}, preprint 2025.

\bibitem{Drinfeld1985}
V.~G.~Drinfeld, \emph{Hopf algebras and the quantum Yang-Baxter
equation}, Dokl.\ Akad.\ Nauk SSSR \textbf{283} (1985), 1060--1064.

\bibitem{STS1983}
M.~A.~Semenov-Tian-Shansky, \emph{What is a classical $r$-matrix?},
Funct.\ Anal.\ Appl.\ \textbf{17} (1983), 259--272.

\bibitem{FF92}
B.~Feigin, E.~Frenkel,
\emph{Affine Kac--Moody algebras at the critical level and
Gelfand--Dikii algebras}, Int.\ J.\ Mod.\ Phys.\ A \textbf{7},
Suppl.\ 1A (1992), 197--215.

\bibitem{FFR1994}
B.~Feigin, E.~Frenkel, N.~Reshetikhin, \emph{Gaudin model, Bethe
ansatz and critical level}, Comm.\ Math.\ Phys.\ \textbf{166}
(1994), 27--62.

\bibitem{Lorgat2026a}
R.~Lorgat, \emph{Modular Koszul duality: a monograph on bar-cobar
duality for chiral algebras on algebraic curves}, Volume~I
(manuscript, 2026).

\bibitem{GZ26}
D.~Gaiotto, P.~Zeng, \emph{Holographic reconstruction of chiral
algebras from commuting sphere Hamiltonians}, preprint 2026.

\end{thebibliography}

\end{document}
